
\section{Rozkład Calderóna--Zygmunda i interpolacja}
\label{sec:dt-calderon-zygmund}
\providecommand{\atlasref}[2]{#2}
Por. \atlasref{PDE-16::lem:pde16-cz-decomp}{RRC-16: rozkład Calderóna--Zygmunda}
oraz \atlasref{ME-25::cor:cubes-differentiation-25}{ME-25: różniczkowanie po kostkach}.
\subsection{Rozkład dyadyczny}
\begin{lemat}[Rozkład Calderóna--Zygmunda]
\label{lem:dt-cz-cz-decomp}
Niech $f\in L^1(\R^d)\cap L^2(\R^d)$ i $s>0$. Istnieje
co najwyżej przeliczalna rodzina parami rozłącznych półotwartych
kostek dyadycznych $Q_k$ oraz rozkład $f=g+\sum_k b_k$
zbieżny w $L^1$ i $L^2$, taki że
\begin{enumerate}[label=(\roman*)]
\item $|g|\le 2^ds$ prawie wszędzie i $\|g\|_1\le\|f\|_1$;
\item $b_k=0$ prawie wszędzie poza $Q_k$,
 $\int_{Q_k} b_k\,\mathrm d\lambda_d=0$,
 $\sum_k\|b_k\|_1\le2\|f\|_1$;
\item $\sum_k|Q_k|\le s^{-1}\|f\|_1$ oraz
 $\|g\|_2^2\le 2^ds\|f\|_1$.
\end{enumerate}
\end{lemat}
\begin{proof}
Ustalamy siatkę kostek $2^m(k+[0,1)^d)$,
gdzie $m\in\mathbb Z$, $k\in\mathbb Z^d$.
Dwie kostki z tej siatki są rozłączne lub jedna
zawiera się w drugiej. Każda ma jedynego rodzica
o dwukrotnie większym boku. Dla $w\in L^1$
i $s>0$ nazywamy kostkę złą, jeśli średnia $|w|$
na niej przekracza $s$. Każda zła kostka zawiera
się w maksymalnej złej kostce: w łańcuchu rodziców
objętości rosną do nieskończoności, a średnie
nie przekraczają $\|w\|_1/|Q|$, więc tylko skończenie
wiele z nich może być złych. Maksymalne złe kostki
są parami rozłączne. Oznaczając przez $M_{\mathrm{dy}}w(x)$
supremum średnich $|w|$ po kostkach dyadycznych
zawierających $x$, dostajemy
\[
 \lambda_d(\{M_{\mathrm{dy}}w>s\})
 =\sum_{Q\ \mathrm{maks.}}|Q|
 \le s^{-1}\sum_{Q\ \mathrm{maks.}}\int_Q|w|\,\mathrm d\lambda_d
 \le s^{-1}\|w\|_1.
\]
Supremum jest mierzalne, gdyż siatka jest przeliczalna.

Uzasadnijmy także różniczkowanie całek po tych kostkach.
Niech $Q_n(x)$ oznacza jedyną kostkę o boku $2^{-n}$
zawierającą $x$ i $A_n w(x)$ średnią $w$ na $Q_n(x)$.
Dla $v\in C_c^\infty(\R^d)$ mamy $A_nv(x)\to v(x)$
w każdym punkcie przez ciągłość i średnicę kostki dążącą
do zera. Dla dowolnego $w\in L^1$ wybieramy
mierzalnego reprezentanta; zachodzi prawie wszędzie
\[
 \limsup_n |A_nw(x)-w(x)|
 \le M_{\mathrm{dy}}(w-v)(x)+|w(x)-v(x)|.
\]
Miara zbioru, gdzie lewa strona przekracza $\tau>0$,
jest więc ograniczona przez $4\|w-v\|_1/\tau$,
na mocy powyższej nierówności i nierówności Czebyszewa.
Z gęstości $C_c^\infty$ w $L^1$
(twierdzenie~\ref{thm:dt-test-density}) miara ta wynosi zero.
Biorąc $\tau=1/l$ dla $l\in\N$, otrzymujemy
$A_nw\to w$ prawie wszędzie.

Wracamy do $w=f$ i oznaczamy maksymalne złe kostki
przez $Q_k$. Ich rodzice nie są źli, więc
\[
 s<\frac1{|Q_k|}\int_{Q_k}|f|\,\mathrm d\lambda_d
 \le 2^ds.
\]
Poza sumą tych kostek wszystkie średnie $|f|$ są
nie większe niż $s$; zastosowane do $|f|$ różniczkowanie
daje $|f|\le s$ prawie wszędzie. Kładziemy
$g=f$ poza kostkami oraz $g=f_{Q_k}$ na $Q_k$, gdzie
$f_{Q_k}=|Q_k|^{-1}\int_{Q_k}f\,\mathrm d\lambda_d$,
a $b_k=(f-f_{Q_k})\chi_{Q_k}$.
Z definicji średniej wynika zerowa całka $b_k$.
Ponadto $|f_{Q_k}|\le 2^ds$ i
\[
 \|g\|_1
 \le\int_{\R^d\setminus\bigcup_k Q_k}|f|\,\mathrm d\lambda_d
    +\sum_k\int_{Q_k}|f|\,\mathrm d\lambda_d
 =\|f\|_1,\qquad
 \|b_k\|_1\le2\int_{Q_k}|f|\,\mathrm d\lambda_d.
\]
Otrzymujemy $\|g\|_2^2\le2^ds\|f\|_1$ oraz
oszacowanie sumy norm $b_k$. Oszacowanie objętości
wynika z nierówności dla $M_{\mathrm{dy}}f$.
Każdy $b_k$ należy do $L^2$, a punktowo prawie wszędzie
$\sum_k b_k=f-g$, przy czym w każdym punkcie co najwyżej
jeden składnik jest niezerowy.
Dla $q=1,2$ z rozłączności i twierdzenia Tonellego
\[
 \sum_k\|b_k\|_q^q=\|f-g\|_q^q<\infty .
\]
Potęga $q$ normy ogona jest sumą odpowiednich potęg
norm, więc szeregi zbiegają w obu przestrzeniach.
Wybierając reprezentanta $b_k$ równego zero poza $Q_k$,
mamy jego nośnik topologiczny zawarty w $\overline{Q_k}$,
niekoniecznie w samej półotwartej kostce.
\end{proof}

\subsection{Oszacowanie słabego typu}
\begin{twierdzenie}[Słaby typ $(1,1)$ operatora Calderóna--Zygmunda]
\label{thm:dt-cz-weak11}
Niech $d\ge1$ i $K\in C^1(\R^d\setminus\{0\};\C)$
spełnia $|K(x)|\le A|x|^{-d}$ i
$|\nabla K(x)|\le A|x|^{-d-1}$, gdzie $A>0$.
Niech $T:L^2(\R^d)\to L^2(\R^d)$ będzie ograniczonym
operatorem liniowym, którego działanie dla każdego
$f\in C_c^\infty(\R^d)$ poza $\supp f$ jest dane przez
zwykłą całkę $\int K(x-y)f(y)\,\mathrm d\lambda_d(y)$
prawie wszędzie.
Istnieje $C>0$, zależna tylko od $d,A,\|T\|_{2\to2}$,
taka, że dla każdego $f\in L^1\cap L^2$ i $t>0$ zachodzi
\begin{equation}\label{eq:dt-cz-weak11}
 \lambda_d\bigl(\{x:|Tf(x)|>t\}\bigr)
 \le \frac{C}{t}\,\|f\|_{L^1(\R^d)}.
\end{equation}
\end{twierdzenie}
\begin{proof}
Stosujemy lemat~\ref{lem:dt-cz-cz-decomp} z $s=t$.
W dowodzie litera $C$ może oznaczać różne dodatnie
stałe zależne wyłącznie od podanych parametrów.
Niech $Q_k^*$ oznacza domkniętą kostkę o tym samym środku
$c_k$ co $Q_k$ i boku $3\sqrt d\,\ell_k$,
gdzie $\ell_k$ jest bokiem $Q_k$. Dla
$E=\bigcup_k Q_k^*$ mamy
$|E|\le(3\sqrt d)^d\sum_k|Q_k|\le C\|f\|_1/t$.
Dla części dobrej nierówność Czebyszewa i
ograniczoność $T$ na $L^2$ dają
\[
 \lambda_d(\{|Tg|>t/2\})
 \le 4t^{-2}\|Tg\|_2^2
 \le Ct^{-1}\|f\|_1.
\]

Uzasadnijmy reprezentację przedłużonego operatora $T$
dla $b_k$, które nie muszą być gładkie.
Wygładzamy $b_k$ standardowymi jądrami
$\rho_\varepsilon$ o nośnikach w $B_\varepsilon(0)$.
Funkcje $b_k*\rho_\varepsilon\in C_c^\infty$
zbiegają do $b_k$ w $L^1$ i $L^2$; wynika to
z twierdzenia~\ref{thm:dt-approx-lp} na ustalonym
zwartym otoczeniu $\overline{Q_k}$, gdyż wszystkie
nośniki są w nim zawarte dla małych $\varepsilon$.
Z ciągłości $T$ na $L^2$ i wyboru podciągu
$T(b_k*\rho_\varepsilon)\to Tb_k$ prawie wszędzie.
Dla ustalonego $x\notin Q_k^*$ odległość od
$\overline{Q_k}$ jest dodatnia; dla małych
$\varepsilon$ jądro $K(x-y)$ jest ograniczone
na wspólnym zwartym otoczeniu nośników.
Zbieżność w $L^1$ daje tam granicę zwykłych całek.
Stąd dla prawie każdego $x\notin Q_k^*$, dzięki zerowej
całce $b_k$,
\[
 Tb_k(x)=\int_{Q_k}[K(x-y)-K(x-c_k)]
 b_k(y)\,\mathrm d\lambda_d(y).
\]
Dla $y\in Q_k$ mamy $|y-c_k|\le\sqrt d\,\ell_k/2$,
a dla $x\notin Q_k^*$ zachodzi
$|x-c_k|\ge3\sqrt d\,\ell_k/2$.
Każdy punkt odcinka od $x-c_k$ do $x-y$ ma więc
moduł co najmniej $2|x-c_k|/3$. Stosując podstawowe
twierdzenie rachunku całkowego wzdłuż odcinka oraz
oszacowanie gradientu jądra, dostajemy
\[
 |Tb_k(x)|\le C\int_{Q_k}
 \frac{|y-c_k|\,|b_k(y)|}{|x-c_k|^{d+1}}
 \,\mathrm d\lambda_d(y)
 \le C\ell_k |x-c_k|^{-d-1}\|b_k\|_1.
\]
Całkując i powiększając zbiór całkowania do
$\{|x-c_k|\ge3\sqrt d\,\ell_k/2\}$, otrzymujemy
\[
 \int_{\R^d\setminus Q_k^*}|Tb_k|\,\mathrm d\lambda_d
 \le C\ell_k\|b_k\|_1
 \int_{3\sqrt d\,\ell_k/2}^\infty r^{-2}\,\mathrm d\lambda_1(r)
 \le C\|b_k\|_1.
\]
Dla skończonych sum nierówność trójkąta daje
$\int_{\R^d\setminus E}|T(\sum_{k\le N}b_k)|
\,\mathrm d\lambda_d\le C\sum_k\|b_k\|_1$.
Ponieważ $\sum_{k\le N}b_k\to f-g$ w $L^2$,
obrazy zbiegają w $L^2$ do $T(f-g)$.
Wybieramy podciąg zbieżny prawie wszędzie
i stosujemy lemat Fatou. Otrzymujemy
\[
 \int_{\R^d\setminus E}|T(f-g)|\,\mathrm d\lambda_d
 \le C\sum_k\|b_k\|_1\le C\|f\|_1.
\]
Z liniowości $Tf=Tg+T(f-g)$ prawie wszędzie.
Zbiór $\{|Tf|>t\}$ zawiera się z dokładnością
do zbioru miary zero w sumie $E$,
$\{|Tg|>t/2\}$ i
$\{x\notin E:|T(f-g)(x)|>t/2\}$.
Ostatnią miarę szacujemy przez $2/t$ razy
poprzednią całkę. To kończy dowód.
\end{proof}

\subsection{Interpolacja i oszacowania w przestrzeniach Lebesgue'a}
\begin{twierdzenie}[Ograniczoność w $L^p$ dla $1<p<\infty$]
\label{thm:dt-cz-general-lp}
Przy hipotezach twierdzenia \ref{thm:dt-cz-weak11}
operator $T$ ma dla każdego $1<p<\infty$ jednoznaczne
ograniczone przedłużenie z $L^2\cap L^p$ na $L^p$,
zgodne z działaniem na $L^2$ na wspólnej dziedzinie.
\end{twierdzenie}
\begin{proof}
Dla $1<p<2$ wynik daje lemat
\ref{lem:dt-cz-interp} zastosowany do oszacowania
\ref{thm:dt-cz-weak11}; dla $p=2$ jest to założenie.
Dla $2<p<\infty$ oznaczmy $q=p/(p-1)$.
Operator sprzężony $T^*$ na $L^2$ jest ograniczony
i spełnia hipotezy twierdzenia \ref{thm:dt-cz-weak11}
z jądrem $K^*(z)=\overline{K(-z)}$.
Istotnie, dla testów $f,h$ o rozłącznych nośnikach
z definicji sprzężenia oraz reprezentacji $Th$ mamy
\[
 \langle T^*f,h\rangle_2
 =\int_{\R^d}\int_{\R^d}f(y)
 \overline{K(y-x)}\,\overline{h(x)}
 \,\mathrm d\lambda_d(y)\,\mathrm d\lambda_d(x).
\]
Twierdzenie Fubiniego ma zastosowanie, gdyż nośniki
są zwarte i oddalone od siebie o dodatnią odległość.
Dla ustalonego $f$ równość zachodzi dla każdego testu
$h$ o nośniku w otwartym dopełnieniu $\supp f$.
Na tym zbiorze funkcja
$x\mapsto\int K^*(x-y)f(y)\,\mathrm d\lambda_d(y)$
jest lokalnie całkowalna, ponieważ na każdym jego
zwartym podzbiorze jądro jest jednostajnie ograniczone
dla $y\in\supp f$. Z jednoznaczności dystrybucji
regularnych (stwierdzenie \ref{prop:regular-injective19})
wynika żądana reprezentacja prawie wszędzie.
Oszacowania $K^*$ i jego gradientu wynikają
bezpośrednio z odpowiednich oszacowań $K$.
Wobec tego $T^*$ jest ograniczony w $L^q$
na mocy pierwszej części dowodu.

Dla $f\in L^2\cap L^p$ oraz $h\in L^2\cap L^q$
otrzymujemy
\[
 |\langle Tf,h\rangle_2|
 =|\langle f,T^*h\rangle_2|
 \le\|f\|_p\|T^*h\|_q
 \le C_q\|f\|_p\|h\|_q.
\]
Aby nie zakładać z góry przynależności $Tf$ do $L^p$,
połóżmy $v=Tf\in L^2$ i
$E_n=B_n(0)\cap\{|v|\le n\}$. Funkcja
$h_n=\chi_{E_n}|v|^{p-2}v$ należy do $L^2\cap L^q$,
gdyż jest ograniczona i ma nośnik skończonej miary.
Dla $a_n=\int_{E_n}|v|^p\,\mathrm d\lambda_d$ mamy
$\langle v,h_n\rangle_2=a_n$ oraz $\|h_n\|_q=a_n^{1/q}$.
Oszacowanie daje $a_n^{1/p}\le C_q\|f\|_p$
(gdy $a_n=0$, nierówność jest automatyczna).
Z twierdzenia o zbieżności monotonicznej dostajemy
$\|Tf\|_p\le C_q\|f\|_p$.
Gęstość $L^2\cap L^p$ w $L^p$ wynika z obcinania
do kul i ograniczonych wartości. Zupełność $L^p$
daje jednoznaczne przedłużenie; jest ono zgodne
na $L^2\cap L^p$ z konstrukcji.
\end{proof}

\begin{uwaga}
Powyższe twierdzenie dotyczy przedłużenia operatora,
który już jest ograniczony w $L^2$ i ma reprezentację
jądrową poza nośnikiem. Nie wynika z niego automatycznie
istnienie granicy głównych wartości prawie wszędzie
dla dowolnej funkcji $L^p$. Do tego potrzebna jest
osobna kontrola obciętych operatorów.
\end{uwaga}

\subsection{Zastosowanie do potencjału Newtona i przypadki końcowe}
\begin{twierdzenie}[Oszacowanie Calderóna--Zygmunda dla laplasjanu]
\label{thm:dt-cz-cz}
Niech $d\ge2$ i $1<p<\infty$. Dla
$f\in C_c^\infty(\R^d)$ oraz $u=\Gamma*f$ zachodzi
\begin{equation}\label{eq:dt-cz-cz}
 \|D^2u\|_{L^p(\R^d)}\le C_{d,p}\|f\|_{L^p(\R^d)}.
\end{equation}
Oszacowanie przedłuża operator $f\mapsto D^2(\Gamma*f)$
do $L^p$; nie oznacza ono, że samo $u$ musi należeć
do $L^p(\R^d)$.
\end{twierdzenie}
\begin{proof}
Jądra $K_{ij}$ i operatory $T_{ij}$ spełniają hipotezy
twierdzenia~\ref{thm:dt-cz-weak11}: oszacowania jądra
wykazaliśmy przy wzorze \eqref{eq:dt-cz-kernel}, reprezentacja
poza nośnikiem wynika z definicji głównej wartości,
a ograniczoność w $L^2$ z lematu~\ref{lem:dt-cz-l2}.
Z ogólnego twierdzenia~\ref{thm:dt-cz-general-lp}
wynika więc, dla każdego $1<p<\infty$,
$\|T_{ij}f\|_p\le C_{d,p}\|f\|_p$.
Stałą możemy wybrać wspólną dla wszystkich $i,j$,
ponieważ dla ustalonego $d$ jest tylko skończenie wiele par.
Ze wzoru~\eqref{eq:dt-cz-repr} mamy
$\|\partial_{ij}u\|_p\le(C_{d,p}+1/d)\|f\|_p$.
Norma Frobeniusa hesjanu jest punktowo nie większa
niż suma modułów jego $d^2$ składowych, więc nierówność
trójkąta w $L^p$ daje~\eqref{eq:dt-cz-cz}.
Gęstość testów (twierdzenie~\ref{thm:dt-test-density})
i zupełność przestrzeni pól macierzowych $L^p$
zapewniają jednoznaczne ograniczone przedłużenie.
\end{proof}

\begin{uwaga}[Dlaczego przypadki $1$ i $\infty$ są osobne?]
\label{rem:dt-cz-endpoints}
Silne oszacowanie $L^1\to L^1$ jest fałszywe nawet
dla gładkich danych. Istotnie, wybierzmy
$f\in C_c^\infty(B_1(0))$, $f\ge0$, $\int f\,\mathrm d\lambda_d=1$.
Dla $|x|>2$ mamy
\[
 |T_{12}f(x)-K_{12}(x)|
 \le C_d|x|^{-d-1}\int |y|f(y)\,\mathrm d\lambda_d(y)
 \le C_d|x|^{-d-1}.
\]
Wynika to z oszacowania gradientu jądra na odcinku
od $x$ do $x-y$, którego punkty mają moduł co najmniej
$|x|/2$. Na pewnym otwartym fragmencie sfery wokół
$(e_1+e_2)/\sqrt2$ mamy $K_{12}(\theta)\ge c_d>0$.
Dzięki jednorodności, dla odpowiednio dużego $r$ i takich
$\theta$, $T_{12}f(r\theta)\ge(c_d/2)r^{-d}$.
Całka po tym stożku jest rozbieżna, ponieważ
$\int_R^\infty r^{-1}\,\mathrm d\lambda_1(r)=\infty$.
Ponieważ $\partial_{12}(\Gamma*f)=T_{12}f$,
nie zachodzi też silne oszacowanie Hessianu w $L^1$.
Oszacowanie~\eqref{eq:dt-cz-weak11} jest słabszym
oszacowaniem końcowym.

Dla końca $p=\infty$ w wymiarze $d=2$ weźmy
$f=\chi_{(0,1)^2}$ i $u=\Gamma*f$.
Poza domkniętym kwadratem wolno różniczkować pod całką
(lokalnie odległość od nośnika jest dodatnia).
Dla $x=(-t,-t)$, $0<t<1/8$, otrzymujemy
\[
 \partial_{12}u(x)=\frac1\pi
 \int_{(0,1)^2}
 \frac{(t+y_1)(t+y_2)}
 {((t+y_1)^2+(t+y_2)^2)^2}\,\mathrm d\lambda_2(y).
\]
Całkowana funkcja jest nieujemna. Na podsektorze
$\pi/6<\theta<\pi/3$, $2t<r<1/2$, gdzie
$y=(r\cos\theta,r\sin\theta)$, obie współrzędne $y$
są co najmniej $r/2$, a $|x-y|\le r+\sqrt2t\le2r$.
Stąd całkowana funkcja jest nie mniejsza niż $c r^{-2}$,
z $c>0$ niezależnym od $t$. Zatem
$\partial_{12}u(-t,-t)\ge c'\log(1/(4t))\to+\infty$.
Pochodna jest ciągła w otwartym dopełnieniu kwadratu,
więc każda dostatecznie duża wartość daje także
przekroczenie zadanego progu na otwartym zbiorze
dodatniej miary. Jest to brak ograniczoności istotnej,
a nie jedynie rozbieżność na ciągu punktów.
Przykład można przenieść na gładkie dane: wygładzenia
$f_\varepsilon=f*\rho_\varepsilon$ mają normę
$L^\infty$ nie większą niż $1$, a w każdym ustalonym
punkcie poza domkniętym kwadratem
$\partial_{12}(\Gamma*f_\varepsilon)$ zbiega
do powyższej całki dzięki zbieżności w $L^1$
i temu, że dla dostatecznie małych $\varepsilon$
wszystkie nośniki leżą w ustalonym zwartym otoczeniu
kwadratu, oddalonym od rozważanego punktu.
Jednostajne oszacowanie $L^\infty$ dla gładkich danych
prowadziłoby zatem do sprzeczności.
\end{uwaga}
